\chapter{Left Ventricular Assist Device (LVAD) Implantation}
\section*{Overview}
An LVAD is an implanted mechanical pump that helps the left ventricle circulate blood to the aorta.
\section*{Common Instruments and Equipment}
The exact setup varies by technique and setting. Common items may include:
\begin{itemize}[nosep,leftmargin=*]
  \item cardiac surgery and cardiopulmonary bypass equipment
  \item LVAD pump, inflow cannula, and outflow graft
  \item controller, driveline, and external power units
\end{itemize}

\section*{Why It Is Done}
It may support patients with advanced heart failure as a bridge to transplant or as long-term destination therapy when transplantation is not available.
\section*{How It Is Performed}
The pump is implanted inside the chest and connected to the left ventricle and aorta. A driveline passes through the skin to external batteries and a controller.
\section*{Contraindications}
No single contraindication applies to every patient; the treating team assesses whether the procedure is suitable for the individual and the reason for treatment.

\section*{Precautions}
The care team reviews relevant health conditions, medicines, allergies, and procedure-specific risks, and explains any changes needed for the individual.

\section*{Risks and Follow-up}
Stroke, bleeding, infection, pump thrombosis, right-heart failure, arrhythmias, and device malfunction are possible. Anticoagulation, driveline care, and specialized follow-up are lifelong.
\section*{References}
\begin{enumerate}\item MedlinePlus. Left Ventricular Assist Device. U.S. National Library of Medicine; 2025.\item Current Medical Diagnosis and Treatment. McGraw-Hill; 2025.\end{enumerate}
\clearpage
